\documentclass[10pt,a4paper]{amsart}
\usepackage[margin=19mm]{geometry}
\usepackage{amsmath,amssymb,mathtools}
\usepackage[scr=rsfs,cal=euler]{mathalfa}
\usepackage{xfrac}
\usepackage[unicode,hidelinks,bookmarks=false]{hyperref}
\newcommand{\ie}{i.e.\ }
\newcommand{\cf}{cf.\ }
\newcommand{\resp}{respectively\ }
\newcommand{\Subsection}{\subsection}
\providecommand{\nobreakdash}{\nobreak\hbox{-}}
\setlength{\emergencystretch}{2em}
\multlinegap0pt
\usepackage{amssymb}
\usepackage{mathtools}
\usepackage{xcolor}
\usepackage{enumerate}
\usepackage{tikz}
\usepackage{float}
\newtheorem{proposition}{Proposition}
\newtheorem{claim}{Claim}
  \newcommand{\go}{{\gothic o}}
  \newcommand{\gothic}{\mathfrak}
  \newcommand{\param}{{\mathchoice{\mkern1mu\mbox{\raise2.2pt\hbox{$
  \centerdot$}}
  \mkern1mu}{\mkern1mu\mbox{\raise2.2pt\hbox{$\centerdot$}}\mkern1mu}{
  \mkern1.5mu\centerdot\mkern1.5mu}{\mkern1.5mu\centerdot\mkern1.5mu}}}
\title{First passage percolation preserves sublinearly Morse boundaries}
\author{Sagnik Jana, Yulan Qing}
\date{Source: arXiv:2603.24423v1 (CC BY 4.0)}
\begin{document}
\maketitle

\noindent\emph{Source and licence.} First passage percolation preserves sublinearly Morse boundaries. Version arXiv:2603.24423v1 (CC BY 4.0), distributed by arXiv under the Creative Commons Attribution 4.0 International licence, http://creativecommons.org/licenses/by/4.0/, as recorded in the arXiv licence metadata for this exact version and shown on the abstract page https://arxiv.org/abs/2603.24423v1. Full licence notice: \texttt{licenses/arxiv-2603.24423-CC-BY-4.0.txt}.\par\smallskip

\noindent\textit{Editorial scope.} Counted unit: the article's claim 313 (source0.tex lines 1485--1487). The article's complete original proof of it is delivered with it (source0.tex lines 1488--1492, one \texttt{proof} environment of 5 lines). No mathematics is written, repaired or re-derived: the statement and the proof are the article's own lines, in the article's own notation, and no retained line is substituted or re-flowed.  Retained uncounted as the source-local context the proof needs: lines 1017--1022.  Nothing was omitted from any retained span, so the package carries no omission record.  No retained line contains a citation, so the wrapper carries no bibliography and no bibliography entry.  No retained line contains a figure environment, an image-inclusion command or drawing code, so no image is delivered and the package declares no asset dependency.  Editorial formatting only, in addition: an AMS \texttt{amsart} preamble that restates the article's own declarations and macros used by the retained lines, namely the environments \texttt{proposition}, \texttt{claim} on separate counters, together with the standard packages those lines need.  The retained environments print in this document as Proposition 1, Claim 1 in order of appearance, whereas the article prints the same items with the numbering of its own full document.  The complete native source of the article as distributed in the arXiv e-print for this exact version is bundled unchanged as \texttt{sources/197144/source0.tex}.
\begin{proposition}\label{extension}\label{BT2}
 Let $X$ be an infinite connected graph with bounded degree, and let $\go$ be some vertex of $X$. Assume $\nu(\{0\})=0$. Then for every constant $D$, there exists $c>0$ such that for a.e. $\omega$, there exists $r_{1}=r_{1}(\omega)$ such that for all finite path $\gamma$ such that $d(\gamma, o) \leq D \cdot \ell(\gamma)$, one has

\[ \ell_\omega(\gamma)\geq c(D) \cdot \ell(\gamma)-r_{1}\]
Furthermore, $c(\param)$ is an decreasing function.
\end{proposition}

   \begin{claim} \label{313}
     For a.e $\omega$ there exist $A=A(\omega), B=B(\omega)$ such that, \[\kappa(d_\omega(\mathfrak{o},p_{i_k})\geq d_\omega(p_{i_k},\gamma^\omega_{\frac 13[s_{i_k},e_{i_k}]}) \geq A \cdot d(\mathfrak{o},p)-B\]  
   \end{claim}

    \begin{proof}
    Since  $\alpha_i{_k} \cap \mathcal{N}_{\kappa'} (\gamma^\omega_{\frac 13[s_{i_k},e_{i_k}]}, m) \neq \emptyset$ then there exist $x$ on $\gamma^\omega$ such that, \[ d_\omega(p_{i_k}, \gamma^\omega)=d_\omega(p_{i_k}, x) \leq m \cdot \kappa(d_\omega(\mathfrak{o},p_{i_k}).\] Suppose $\sigma^\omega$ be the geodesic joining $p_{i_k}$ and $x$ in $X_\omega$ and under FPP map the preimage $\sigma$ be the path joining $p_{i_k}$,$x$ in $X$. Since $p_{i_k} \in \alpha_{i_k}$ then $p_{i_k}$ is outside the $c_1$ linear neighborhood. Therefore we get, 
    \[|\sigma|\geq d(p_{i_k},x)\geq d(p_{i_k},\gamma)\geq c_1 \cdot d(\mathfrak{o},p_{i_k}) \geq c_1 \cdot d(\mathfrak{o}, \sigma)\]
    Thus, by proposition ~\ref{BT2} for a.e $\omega$ there exist $r_1=r_1(\omega)$ such that, \[ d_\omega(p_{i_k},x)\geq c \cdot d(p_{i_k},x)-r_1.\] Hence, the claim follows.
    \end{proof}

\end{document}
